% ===== §3 =====
\section{Transmission and Collective Representation}
\label{sec:transmission}

\noindent
The first tradition is the one that gave the modern concept its founding definition, and it is placed first because it holds culture most nearly as a settled thing. When Tylor opened \emph{Primitive Culture} in 1871 with the sentence that culture is ``that complex whole which includes knowledge, belief, art, morals, law, custom, and any other capabilities and habits acquired by man as a member of society'' \citep{tylor1871}, he did three things at once, and each has proved durable. He made culture a \emph{whole}, an ensemble rather than a list, so that its parts are to be understood in their belonging together. He made it \emph{acquired}, learned rather than inherited in the blood, and so severed it from race and biology, which was the definition's great emancipatory achievement and remains its permanent contribution. And he made it acquired \emph{by man as a member of society}, so that the individual receives culture from a collectivity that precedes him. Everything the definition gets right, this paper keeps. What it fixes, and fixes so naturally that the fixing is invisible, is the mode of being of the whole so defined. Culture, in this founding image, is something one \emph{acquires}, and what can be acquired is a thing that was there to be had, standing ready before the acquiring, the same for the next person who will acquire it in turn. The verb quietly installs an object. The complex whole becomes a patrimony, a stock of acquirable items, and the work of the concept from that point forward is the work of transmission: how the stock passes, how it is kept, how it changes across the passing. Culture is modelled on inheritance, and an inheritance is by definition a thing that exists before and after the one who receives it.

Durkheim gave this image its deepest sociological form, and in doing so raised the collective from the background of Tylor's definition into the explicit ground of culture's existence. The concepts, categories, and beliefs by which the members of a society understand their world are, for Durkheim, \emph{collective representations}: they are not the constructions of individual minds but products of the collectivity, which express, and in expressing sustain, the reality of the group itself \citep{durkheim1912}. Their mark is that they confront the individual from without and with a force he did not author, the defining property Durkheim assigned to the social fact, which is exterior to any given individual and endowed with a power of constraint \citep{durkheim1982rules}. This is a profound advance over Tylor, because it grasps that culture is not merely deposited in individuals but obtains at the level of the collectivity and bears down on the individual with the weight of an independent reality. It moves culture, in the terms of the previous section's map, from the purely individual toward the collective bearer. And yet it moves it only that far, and in one crucial respect it deepens the very fixity it inherits. For the collective representation, precisely by being exterior and constraining, is figured as a thing standing over against the individual, a structure he enters and by which he is formed, already there before him and outlasting him. The advance in the bearer is paid for by an entrenchment in the mode of being. Where Tylor's culture was an acquirable stock, Durkheim's is an exterior and coercive order, and both are, in the decisive sense, \emph{already made}. The individual meets culture as one meets a fact: it is given, it constrains, and its existence does not wait on his participation.

\subsection*{What this tradition sees, and what its presupposition costs}

It would be an error, and the kind of error this paper is at pains to avoid, to treat this tradition as simply mistaken. It sees something real, and the seeing must be preserved. Culture \emph{is} transmissible; it \emph{does} sediment into durable forms that outlast particular occasions; it \emph{does} confront the individual with the force of something he did not make and cannot simply revise at will. A private language between two people, once established, bears down on each of them with a version of exactly this exteriority: one cannot unilaterally decide that a word the two of you made together now means something else. The transmissibility and the sedimented durability that this tradition places at the centre are genuine features of culture, and any account that denied them, in the name of process and generation, would be a sentimentalism deserving the correction this tradition would administer to it. \emph{Wen}---pattern, the durable mark---is not an illusion.

But the presupposition beneath the seeing exacts a price, and the price is the disappearance of the very thing this paper means to recover. By modelling culture on the acquirable stock and the exterior fact, the tradition renders unthinkable the moment of \emph{emergence}: the event in which a culture that did not exist comes to be, between people who are making it, before it has hardened into anything transmissible or exterior at all. If culture is what is already there to be acquired, then the coming-to-be of the not-yet-there falls outside the concept. One can ask how a culture is transmitted, maintained, or eroded; one cannot, within this frame, ask how it is \emph{generated}, because generation is precisely the passage from the not-yet-cultural to the cultural, and a theory that begins from the already-made has placed that passage before its own first page. And it is the moment of emergence, not the fact of transmission, that the intimate dyad exhibits with unique clarity and that this paper takes as its subject. Two people do not, in the first instance, \emph{acquire} their shared culture; they \emph{make} it, and one can watch them make it. The founding tradition, for all it correctly sees, cannot look at that making, because its concept of culture is built from the residue the making leaves behind. This is the sense in which it is the antithesis's main point of attack: not that transmission is unreal, but that transmission is the afterlife of an emergence the tradition cannot think, and that to take the afterlife for the whole is to mistake the ash for the fire.
